\section{Conclusion and Future Work}
\label{sec:conclusion}
This paper presented the Blind Cognitive Multi-Factor Authentication (BC-MFA) architecture for untrusted web and mobile public clients exposed to unprivileged user-space malware. By decoupling an ephemeral $4 \times 4$ quadrant-balanced toroidal matrix ($H_{\text{bal}} = 34.19\text{ bits}$) from a private quadrant anchor and non-zero offset trajectory, BC-MFA mitigates resident keyloggers, screen scrapers, and carrier-level SIM swapping without requiring dedicated client hardware tokens. A dual-channel PRISMA 2020 review synthesized 40 authoritative sources (30 peer-reviewed domain studies, 9 normative specifications, and the PRISMA reporting guideline). We formalized a standards-compliant OpenID Connect Core 1.0 public-client profile combining RFC 7636 PKCE, registered RFC 8176 \texttt{amr} claims (\texttt{mfa}, \texttt{kba}, \texttt{rba}, \texttt{swk}), and RFC 9449 DPoP access token binding. Analytical threat modeling and a reproducible 500-trial rational-adversary Monte Carlo simulation (\texttt{seed=2026}) quantify multi-session intersection dynamics, establishing $k_{\max} = 2$ (1.6\% exact recovery; 0.0\% at $k_{\max} = 1$) as the operational re-enrollment boundary.

Three explicit boundaries qualify this design proposal. First, deterministic phishing resistance on first-contact unpinned kiosks remains outside software-only DPoP scope when active proxies rewrite client code. Second, 4-character input limits passive behavioral telemetry primarily to human-vs-automation and relay anomaly detection rather than standalone cross-device biometric user identification. Third, spatial modular arithmetic imposes onboarding and accessibility trade-offs requiring WCAG-compliant alternative assurance pathways. Future work within the Master's dissertation at ISPGAYA will deploy the Keycloak 24 and TypeScript/React Native testbed, executing the IRB- and GDPR-compliant longitudinal study ($N = 64$, within-subjects vs. TOTP/WebAuthn) and live Evilginx/infostealer ablation trials.
